% Part: many-valued-logic
% Chapter: three-valued-logics
% Section: lukasiewicz

\documentclass[../../../include/open-logic-section]{subfiles}

\begin{document}

\olfileid{mvl}{inf}{luk}

\olsection{\L ukasiewicz यांचे तर्कशास्त्र}

\begin{editorial}
  अनंतमूल्य \L ukasiewicz तर्कशास्त्रावरील हा विभाग
  संक्षिप्त ``प्रारूप'' आहे.
\end{editorial}

\begin{defn}\ollabel{def:lukasiewicz} अनंतमूल्य \L ukasiewicz
तर्कशास्त्र~$\LogLuk[\infty]$ पुढील मॅट्रिक्सने परिभाषित होते:
\begin{enumerate}
  \item $\lnot$, $\land$, $\lor$, $\lif$ असलेली मानक विधानीय भाषा $\Lang L_0$.
  \item सत्यतामूल्यांचा संच $V_\infty$.
  \item $1$ हेच एकमेव निर्दिष्ट मूल्य आहे, म्हणजे $V^+ = \{1\}$.
  \item सत्यता-फलने पुढील सूत्रांनी दिली आहेत:
  \begin{align*}
    \tf{\lnot}[\LogLuk](x) & = 1 - x\\
    \tf{\land}[\LogLuk](x,y) & = \min(x,y)\\
    \tf{\lor}[\LogLuk](x,y) & = \max(x,y)\\
    \tf{\lif}[\LogLuk](x,y) & = \min(1,1-(x-y)) = \begin{cases}
      1 & \text{जर } x \le y\\
      1-(x-y) & \text{इतर वेळी.}
    \end{cases}
    \end{align*}
\end{enumerate}
$m$-मूल्य \L ukasiewicz तर्कशास्त्राची व्याख्या याचप्रमाणे,
फक्त $V = V_m$ घेऊन केली जाते.
\end{defn}

\begin{prop}
  \olref[thr][luk]{def:lukasiewicz} मध्ये परिभाषित
  $\LogLuk[3]$ आणि \olref{def:lukasiewicz} मध्ये
  परिभाषित $\LogLuk[3]$ हे एकच तर्कशास्त्र आहे.
\end{prop}

\begin{proof}
  \olref[thr][luk]{def:lukasiewicz} मधील संयोजकांची
  सत्यताकोष्टके आणि \olref{def:lukasiewicz} मधील
  समीकरणांनी निश्चित होणारी सत्यताकोष्टके यांची तुलना
  केल्यास हे दिसते:
  \begin{center}
    \begin{tabular}{c|c}
      $\tf{\lnot}$ & \\
      \hline
      $1$ & $0$ \\
      $1/2$ & $1/2$ \\
      $0$ & $1$
    \end{tabular}
    \quad
    \begin{tabular}{c|ccc}
      $\tf{\land}[\LogLuk[3]]$ & $1$ & $1/2$ & $0$ \\
      \hline
      $1$ & $1$ & $1/2$ & $0$ \\
      $1/2$ & $1/2$ & $1/2$ & $0$\\
      $0$ & $0$ & $0$ & $0$
    \end{tabular}
    \\[2ex]
    \begin{tabular}{c|ccc}
      $\tf{\lor}[\LogLuk[3]]$ & $1$ & $1/2$ & $0$ \\
      \hline
      $1$ & $1$ & $1$ & $1$ \\
      $1/2$ & $1$ & $1/2$ & $1/2$ \\
      $0$ & $1$ & $1/2$ & $0$
    \end{tabular}
    \quad
    \begin{tabular}{c|ccc}
      $\tf{\lif}[\LogLuk[3]]$ & $1$ & $1/2$ & $0$ \\
      \hline
      $1$ & $1$ & $1/2$ & $0$ \\
      $1/2$ & $1$ & $1$ & $1/2$  \\
      $0$ & $1$ & $1$ & $1$
    \end{tabular}
  \end{center}
\end{proof}

\begin{prop}\ollabel{prop:luk-infty-m}
  $\Gamma \Entails[\LogLuk[\infty]] !B$ असेल, तर
  सर्व~$m \ge 2$ साठी $\Gamma \Entails[\LogLuk[m]] !B$.
\end{prop}

\begin{proof}
  सराव.
\end{proof}

\begin{prob}
  \olref[mvl][inf][luk]{prop:luk-infty-m} सिद्ध करा.
\end{prob}

खरे तर आधारसूत्रांचा संच ससीम असेल, तर याचे प्रतिलोम
विधानही खरे आहे.

अनंतमूल्य \L ukasiewicz तर्कशास्त्र हे सर्वाधिक प्रचलित
अस्पष्ट तर्कशास्त्र आहे. अस्पष्ट तर्कशास्त्रावरील साहित्यात
सशर्त विधानाची व्याख्या अनेकदा $\lnot !A \lor !B$ अशी
केली जाते. त्यातून अनंतमूल्य प्रबळ क्लिनी तर्कशास्त्र मिळेल.

\begin{prob}
  $(p \lif q) \lor (q \lif p)$ हे $\LogLuk[\infty]$
  मध्ये सर्वतः सत्य सूत्र आहे हे दाखवा.
\end{prob}

\end{document}
